\Qhead{2.1}{Плоские волны. Поток энергии. Поляризация}{Плоские однородные и неоднородные волны. Плоские акустические волны в вязкой теплопроводящей среде, упругие продольные и поперечные волны в твердом теле, электромагнитные волны в среде с проводимостью. Поток энергии. Поляризация.}
\Subfile{2.1.1}{sections/02_waves/q01/s01}
\Subfile{2.1.2}{sections/02_waves/q01/s02}
\Subfile{2.1.3}{sections/02_waves/q01/s03}
\Subfile{2.1.4}{sections/02_waves/q01/s04}

\recap{Плоская волна определяется геометрией фазы; затухание задаётся мнимой частью волнового вектора. Уравнения среды определяют поляризации и дисперсию. Энергетический поток следует из локального баланса, а не всегда направлен по волновому вектору.}
\board{Комплексный $\bm k$; линеаризованные уравнения жидкости; матрица Христоффеля; $k^2=\omega^2\mu\tilde\varepsilon$ и скин-слой.}
\oral{Почему в жидкости нет обычной поперечной упругой волны, а в твёрдом теле есть?}
\bait{Можно оставить различие однородной и незатухающей волн: комплексный волновой вектор сам по себе не делает волну неоднородной.}
\sources{Цысарь~\cite{tsysar1}, PDF-с.~5--33; проводящая среда: §6.4, с.~32--33; Белокопытов~\cite{belokopytov1}, разделы о волнах. Теплопроводная поправка и уравнение Христоффеля выписаны как стандартные дополнения с явными определениями коэффициентов.}
